% !TeX root = ../main.tex
\section{Beats, bars and rhythm notation}
\label{sec:rhythm}

\subsection{What the stage must do}

A tab needs bars, a time signature per bar and a written duration per note.
That takes three steps: find the \term{beats} (the pulse a listener taps along
to), find the \term{downbeats} (the first beat of each bar, which fixes the
meter: 4 beats per bar for 4/4, 3 for 3/4), and \term{quantise} each note onto
a grid of beat subdivisions so it can be written as a quarter, eighth,
sixteenth or triplet.

\subsection{How beat trackers work}

\begin{keyidea}[: activation, then decoding]
A tracker first computes, for every short frame, how likely a beat is there (an
\term{activation}): classically from an onset function such as spectral flux,
today from a neural network. It then \emph{decodes} a regular sequence of beats
from the activation with dynamic programming (Section~\ref{sec:viterbi}), using
the prior knowledge that beats are roughly evenly spaced.
\end{keyidea}

Ellis's tracker~\cite{ellis2007beat} picks the tempo by autocorrelation of the
onset function and then the beat times by dynamic programming. Neural trackers
predict beat and downbeat activations and decode them with a dynamic Bayesian
network that assumes a fixed meter and a tempo range~\cite{bock2016joint};
All-In-One~\cite{kim2023allinone} also tracks song sections.
\textbf{Beat This!}~\cite{foscarin2024beatthis} drops the DBN: a transformer
trained on diverse data (including solo instruments and time-signature changes)
predicts beats and downbeats that are only peak-picked, and surpasses prior
work in F1 while being weaker on continuity \reported.
AlignBeat~\cite{ahn2025alignbeat} reports 3.6 points more downbeat CMLt than
Beat This! \reported; Explore This!~\cite{kihlborg2026explore} predicts a finer
\term{tatum} grid and groups it into beats and bars. A failure analysis on the
SMC dataset~\cite{ahn2026smc} found that current trackers are ``confident but
wrong'' on expressive music, and that the usual DBN minimum tempo of 55\,BPM
forces double-tempo output on 21\% of its tracks \reported.

\subsection{From beats to written rhythm}

Quantisation and the choice of notated durations are rarely studied in the
tablature literature; TART uses a rule-based tab generator~\cite{gupta2026tart}.
The difficulty is that every beat can be divided in two ways that matter in
popular music: \term{straight} (two eighths, four sixteenths) or
\term{triplet} (three eighth-note triplets). A note 1/3 of a beat late is a
triplet in a shuffle and a timing error in a straight rock song.

\begin{inapp}[: rhythm \done]
\begin{itemize}
\item \textbf{Beat This!} (\file{transcribe/beat_net.rs}; ONNX files
  \code{beat_this_final0.onnx}, or \code{beat_this_small0.onnx} when final0 is
  missing, plus \code{beat_this_logmel.onnx}): runs on the mix at 22.05\kHz{},
  128 log-mel bands, 50 activations per second, in 30\,s chunks. Peaks become
  beats; \code{regularize} fills missed beats and drops doubled ones.
\item \textbf{Bars and meter} (\code{beat_net::decode_downbeats}): a
  \term{bar-pointer} Viterbi whose states are ``beat $k$ of a bar of $n$
  beats'', $n$ from 2 to 7. Each meter has a cost (4/4 is favoured,
  \code{METERS}) and changing meter costs \code{METER_CHANGE}; downbeat
  activations are scored against their mean over the song. So a song can switch
  meter or insert a single 2/4 bar.
\item \textbf{Fallback} (\file{transcribe/beats.rs}) when Beat This! is not
  installed: spectral flux, autocorrelation tempo with a 120\,BPM prior, Ellis
  dynamic programming, downbeat phase from the low-frequency band; 4/4 only. It
  runs on the drum stem when there is one.
\item \textbf{Grid} (\code{beats::bar_grid}, \code{arrange::BeatGrid}): whole bars
  around the notes, one length per bar.
\item \textbf{Quantisation} (\file{transcribe/arrange.rs}): onsets within 35\ms{}
  of each other form a chord; \code{choose_subdivisions} picks, for each beat,
  sixteenths or eighth-note triplets by a Viterbi search with a prior towards
  straight time; positions are stored in ticks (960 per beat). \code{notate}
  writes tuplets only inside partly used triplet beats and ties long notes
  across beats; each bar gets its own tempo in the score.
\end{itemize}
\end{inapp}

\subsection{What we measured}

\begin{table}[h]
\centering
\caption{Beat tracking on GuitarSet players 03--05, all 180 excerpts
\measured{} (\code{bench_beats.py}, \code{mir_eval}).}
\label{tab:bench-beats}
\small
\begin{tabular}{@{}lcc@{}}
\toprule
 & Accompaniment (comp) & Solo \\
Tracker & beat F / CMLt / downbeat F & beat F / CMLt / downbeat F \\
\midrule
Beat This! small0 & 0.92 / 0.80 / 0.92 & 0.51 / 0.36 / 0.33 \\
Beat This! final0 (app default) & \textbf{0.97 / 0.92 / 0.94} & \textbf{0.55} / 0.35 / \textbf{0.36} \\
Ellis DP (librosa) & 0.54 / 0.42 / -- & 0.39 / 0.26 / -- \\
\bottomrule
\end{tabular}
\end{table}

\begin{itemize}
\item final0 is a clear gain on accompaniment (CMLt 0.80 to 0.92) for 8\,MB more
  weights.
\item Every tracker fails on unaccompanied solo lines (beat F about 0.5): a
  single melody carries little beat information. In a song the tracker runs on
  the full mix, so this matters mostly for solo-guitar recordings.
\item In the app pipeline (\code{eval_guitarset}, players 03--05, 60 excerpts,
  GAPS), final0 against small0: tempo right in 51 against 49 excerpts, bar lines
  aligned in 41 against 35 (solo 11 against 6, comp 30 against 29) \measured.
  The beat stage takes 13.5\,s against 8.3\,s on a 261\,s song.
\end{itemize}

\begin{caution}[: two sets of rhythm numbers]
The README and an older ROADMAP table report ``tempo right in 59 of 60, bar
lines aligned in 56--57 of 60'' on GuitarSet. Those came from an earlier run
with a different selection of 60 excerpts (and small0); the numbers above are
on players 03--05 only. Even small0 scores only 49 of 60 on players 03--05, so
the selection, not the model, explains most of the gap, but nobody has checked
which excerpts differ. Quote the players 03--05 numbers as current.
\end{caution}

\subsection{Failure modes}
\begin{itemize}
\item Compound meters (6/8, 12/8) are written as 2/4 or 6/4 with triplets
  \measured.
\item Sixteenth-note triplets and 32nd notes cannot be written (the grid has
  only sixteenths and eighth triplets).
\item Slow, rubato or expressive music: confident but wrong beats, double
  tempo~\cite{ahn2026smc} \reported.
\item Solo-guitar recordings without accompaniment: about half the beats
  missed (Table~\ref{tab:bench-beats}).
\end{itemize}
